% META: {"id":"1SPE-DERLOCAL-EX-051","chapitre":"1SPE-DERIVATION-LOCAL","type_objet":"exercice","capacites":["1SPE-DERIVATION-LOCAL-C3","1SPE-DERIVATION-LOCAL-C4"],"capacites_codes":["C3","C4"],"methodes":[],"parcours":1,"competences":["representer","raisonner"],"duree_min":12,"mode_creation":"ex_nihilo","sources_inspiration":[],"fichier_tex":"chapitres/1SPE-DERIVATION-LOCAL/exercices/1SPE-DERLOCAL-EX-051.tex","corrige_tex":"chapitres/1SPE-DERIVATION-LOCAL/corriges/1SPE-DERLOCAL-CO-051.tex","status":"generated"}
% BEGIN-VERIFY
% # Automatisme : passer du graphique aux donnees et reciproquement.
% # Lecture : la tangente en A(2 ; 3) passe par B(0 ; 7).
% pente = (3 - 7) / (2 - 0)
% assert pente == -2.0
% # Equation de la tangente : y = f'(a)(x - a) + f(a).
% tangente = lambda x: pente * (x - 2) + 3
% assert tangente(2) == 3
% assert tangente(0) == 7
% assert all(abs(tangente(x) - (-2 * x + 7)) < 1e-12 for x in (-3, 0, 2, 5))
% # Retour au graphique : ordonnee a l'origine 7, coefficient directeur -2.
% assert tangente(0) == 7 and tangente(1) - tangente(0) == pente
% # Approximation affine au voisinage de 2.
% assert abs(tangente(2.1) - 2.8) < 1e-12
% assert abs(tangente(1.9) - 3.2) < 1e-12
% # Une fonction possible : f(x) = -x^2 + 2x + 3 verifie les deux lectures.
% f = lambda x: -x * x + 2 * x + 3
% fprime = lambda x: -2 * x + 2
% assert f(2) == 3
% assert fprime(2) == -2
% END-VERIFY
\begin{exercice}{1SPE-DERLOCAL-EX-051}{1}{12}
Sur un graphique, on lit la courbe $\mathcal{C}_f$ d'une fonction $f$
dérivable, et sa tangente $T$ au point $A$ d'abscisse $2$. On relève deux
points de cette tangente : $A(2\,;\,3)$ et $B(0\,;\,7)$.
\begin{enumerate}
\item \textbf{Du graphique aux données.} Déterminer $f(2)$ et $f'(2)$.
\item En déduire l'équation réduite de la tangente $T$.
\item \textbf{Des données au graphique.} On veut retrouver $T$ sans calcul :
  quelle ordonnée à l'origine et quel coefficient directeur faut-il utiliser
  pour la tracer ?
\item Donner une valeur approchée de $f(2{,}1)$ et de $f(1{,}9)$ à l'aide de
  l'approximation affine.
\item Vérifier que la fonction $f(x) = -x^2 + 2x + 3$ est compatible avec les
  deux lectures de la question 1.
\end{enumerate}
\end{exercice}
